\documentclass{article}
\usepackage[preprint]{tmlr}
\usepackage{amsmath,amssymb,amsfonts}
\usepackage{booktabs}
\usepackage{graphicx}
\usepackage{xcolor}
\usepackage{hyperref}
\usepackage{url}
\usepackage{microtype}
\usepackage{enumitem}

% ---- draft macros -------------------------------------------------------
\newcommand{\todo}[1]{{\color{red}\textbf{[TODO: #1]}}}
\newcommand{\planned}[1]{{\color{blue!70!black}\textit{#1}}}
\newcommand{\dmax}{\Delta_{\max}}
\newcommand{\qdmax}{\dot q_{\max}}
\newcommand{\iface}{\boldsymbol{\theta}}

\title{Where Does Manipulation Precision Come From?\\
\large Attributing the Precision Floor of Reinforcement-Learned Policies
to the Action Interface}

\author{\name \todo{Author name} \email \todo{email}\\
\addr Independent Researcher, Tashkent, Uzbekistan}

\def\month{MM}
\def\year{YYYY}
\def\openreview{\url{https://openreview.net/forum?id=XXXX}}

\begin{document}
\maketitle

% ------------------------------------------------------------------------
\begin{center}
\fbox{\parbox{0.95\linewidth}{\small\textbf{Draft status (8 Sep 2026).}
Every number printed in black is regenerated from a per-episode CSV committed
in the public repository by \texttt{python reproduce.py}; none is transcribed
from a console log. The eval-only part of Phase~0
(Sec.~\ref{sec:phase0}) is \emph{complete} --- 51 archived cells, 63\,000
evaluation episodes, no training --- and it is what the released checkpoints
alone can say. Phase~1, Phase~2, the second environment and the retrained solver
extremes require 73 training runs that have not been run; everything describing them is
marked \planned{blue}. \todo{red} marks a decision that is still open. The
August 2026 figures that appeared in earlier versions of this draft are
superseded; \texttt{CHANGELOG.md} in the repository lists every correction.}}
\end{center}

\begin{abstract}
Reinforcement learning can now be pushed to millimetre precision on real
robots: \citet{luo2025hilserl} reach near-perfect success on connector
insertion, timing-belt assembly and dual-arm coordination within one to two
and a half hours of on-robot training, by putting a human in the loop and
choosing the controller interface carefully. What is not known is which of
those ingredients buys the precision. \citet{xu2026curse} give the question a
name, formalising the saturation point as a limit precision $c$ in a
data-scaling law for behaviour cloning and describing $c$ as ``an emergent
property of the entire agent system,'' but leaving open both what determines
it and whether it survives online reinforcement learning. We study $c$ for a reinforcement-learned pick-and-place policy on a
6-DoF arm in MuJoCo and test a single hypothesis: \emph{the floor is set by the
action interface} --- the maximum per-step displacement, the joint-target
smoothing constant, the control rate, and the joint-velocity limit imposed by
the safety layer --- rather than by the reward, the training budget, the
observation encoder, or the contact-solver parameters of the simulator. We
retrain policies at each interface setting instead of relabelling one policy's
errors, evaluate every cell with $n{=}2000$ episodes under fixed domain
randomisation, run four control arms designed to falsify the interface
hypothesis, and repeat the step-size axis in a second, independently authored
environment --- \texttt{gym-hil}, the open simulation companion to the
human-in-the-loop line of work, with a different arm, a different controller
and different interface constants --- so that a result about interfaces is not
a result about one robot. This version reports the complete eval-only phase. Two released
policies with different training histories place the object with median
terminal errors of 17.9 and 15.7\,mm, roughly $0.5$--$0.6$ times the 30\,mm
maximum step; the terminal error is almost entirely the tool-centre-point
error at the instant the object leaves the fingers ($r^2$ up to $0.28$) and
essentially unrelated to how far the object then rolls ($r^2 \approx 0.004$),
so the floor is set at release rather than by post-release settling. The
contact-solver kill gate --- the released policy re-evaluated under integration
timestep, contact time constant, solver iteration count and impedance ratio
changes spanning the range used by published MuJoCo manipulation work --- moves
$c_{50}$ by \input{generated/tables/inline_gate_span}\unskip, while changing how
often the policy fails outright a good deal more. Evaluating all 38 stored
checkpoints shows the floor is undefined until the task is solved, drops once,
and then oscillates without improving over the last third of training; Xu et
al.'s scaling form does not fit an RL checkpoint curve ($r^2 = 0.17$) and we
report the failure rather than a fitted constant. A nested design of 200 scenes
replayed under 10 noise streams each puts the intra-class correlation of the
outcome at the floor at $0.175$: five parts noise to one part scene.
\planned{The interface sweep, the estimated $c(\iface)$ and the training-based
control arms are pending; see Sec.~\ref{sec:results}.} We release the environment, the trained
checkpoints, per-episode evaluation logs with the sampled randomisation
parameters, and a one-command reproduction script.
\end{abstract}

% ------------------------------------------------------------------------
\section{Introduction}
\label{sec:intro}

\citet{luo2025hilserl} showed that reinforcement learning can be driven to
millimetre precision on real hardware --- connector insertion, timing-belt
assembly, dual-arm handovers --- in one to two and a half hours of on-robot
training. Their system changes several things at once: a human intervenes
during learning, a reward classifier is trained from teleoperated samples, and
the policy commands the arm through a carefully chosen impedance interface.
Which of those is doing the work for \emph{precision} specifically is not
separable from that result, and it is not obvious. This paper isolates one
candidate and measures it.

Start with the phenomenon. A learned manipulation policy that succeeds on
94.9\% of pick-and-place episodes when ``success'' means landing within 45\,mm
of the target succeeds on 20.8\% of them when the tolerance is 10\,mm and on
5.3\% when it is 5\,mm ($n{=}2000$; Table~\ref{tab:baselines}). Nothing about
the policy changed; only the ruler did. Every practitioner who has tightened a
tolerance has seen this cliff, and yet the community's benchmarks report a
single success rate at a single, usually generous, tolerance
\citep{jiang2026benchmarking}, which hides the cliff entirely.

\citet{xu2026curse} were the first to characterise the cliff quantitatively.
Training behaviour-cloning policies on three simulated tasks, they show that
the number of demonstrations $N$ needed to reach a precision $P$ diverges as
$\log N \propto 1/(P - c)$: there is a limit precision $c$ that no amount of
data crosses. They measure $c \approx 2.35$\,mm for peg insertion and
$2.75$\,mm for cuboid stacking, show that $c$ moves when the wrist camera is
removed ($\to 3.85$\,mm), when the demonstrator is cleaner ($\to 1.27$\,mm) or
when randomisation is reduced ($\to 1.00$\,mm), and conclude that $c$ ``is not
a static physical constant of the task but an emergent property of the entire
agent system.'' They fix the action space to a 7-DoF end-effector delta pose
throughout, study behaviour cloning only, and name reinforcement learning as
the open next step. They release neither code nor data, so the RL case cannot
be answered by re-analysing theirs; it has to be measured.

This paper asks the question their result leaves open, for reinforcement
learning: \textbf{what component of the agent system sets $c$?} We advance one
falsifiable answer. In a delta-action interface the policy never commands a
position; it commands a bounded displacement $a \in [-1,1]^d$ that is scaled by
a maximum step $\dmax$, filtered by a smoothing constant $\alpha$, rate-limited
by a joint-velocity bound $\qdmax$ and applied at $f$\,Hz. These four numbers
--- the \emph{action interface} $\iface = (\dmax, \alpha, f, \qdmax)$ --- define
the finest correction the arm can make in the final steps before release. Our
hypothesis is that $c$ is a function of $\iface$ and, to first order, of nothing
else: not of the reward, not of the training budget, not of the observation
encoder and not of the simulator's contact parameters.

The hypothesis is worth stating because the obvious alternatives are all live.
\citet{gyenes2026fourier} attribute the precision floor of imitation policies to
the spectral bias of MLP encoders and remove it with Fourier features --- a
representation explanation. \citet{bronars2026tune} show that reinforcement
learning is agnostic to controller gains once hyperparameters are re-tuned ---
suggesting the learner can absorb interface changes, which would predict a flat
$c(\iface)$. A precision term in the reward, or a longer training run, is what
most practitioners try first. And any sim-only claim invites the objection that
the floor is an artefact of contact stiffness in the MJCF. Each of these is a
control arm in our design (Sec.~\ref{sec:design}), and the last of them ---
the one that would end the study --- is reported here in full.

Three features of our setup make the question answerable on one workstation.
First, the policy is trained by soft actor-critic \citep{haarnoja2018sac} on a
state observation, so a full retraining costs hours rather than days and we can
afford to \emph{retrain} at every interface cell rather than relabel one
policy's errors post hoc --- the latter cannot support any claim about training
budget. Second, the environment records the sampled domain-randomisation vector
of every episode, so each success curve is accompanied by a failure
decomposition over the randomisation space (Sec.~\ref{sec:attrib}). Third, the
step-size axis is repeated in a second environment we did not write:
\texttt{gym-hil} \citep{gymhil2025}, the open MuJoCo suite that accompanies the
human-in-the-loop line of work, whose Franka Panda runs under operational-space
control with a mocap target rather than a Piper under differential IK. Its
interface constants were chosen independently of ours, and
Sec.~\ref{sec:twoenvs} reports what they are.

That comparison is worth stating before any result. The Piper environment of
this paper commands at most 30\,mm per step at 20\,Hz with a smoothing constant
of 0.45, and calls a placement successful within 45\,mm. \texttt{gym-hil}
commands at most 25\,mm per step at 10\,Hz with no smoothing at all, and calls
a grasp successful within 50\,mm and an arrangement within 30\,mm. Two
environments, two authors, two arms, two controllers --- and every success
tolerance in both lands between one and two commanded step sizes. If the
precision floor were set by the reward, the budget or the simulator, there
would be no reason for the people who wrote these environments to have
converged on that ratio.

\paragraph{Contributions.}
\begin{enumerate}[nosep,leftmargin=*]
\item \textbf{(Measured.)} A complete eval-only characterisation of the floor of
two released policies at $n{=}2000$: success-versus-tolerance curves with
Wilson intervals over $\tau \in [5, 45]$\,mm, $c_{50}$ and $c_{90}$ with
bootstrap intervals, a decomposition of the terminal error into
tool-centre-point error at release and post-release scatter, and a
precision-versus-training-steps curve over all 38 stored checkpoints of both
runs, and an intra-class correlation of the outcome from a nested
scene-by-noise design.
\item \textbf{(Measured.)} A pre-registered contact-solver kill gate: the
released policy re-evaluated under integration timestep, contact time constant,
solver iteration and impedance-ratio changes, testing whether the floor is a
simulator artefact before any training budget is spent on the main sweep.
\item \planned{(Pending, 43 training runs.)} A retraining-based measurement of
$c$ across a one-factor-at-a-time sweep of the action interface, with the three
remaining control arms --- reward shaping, training budget and Fourier-feature
encoding --- and a hyperparameter-optimisation arm at the smallest step size.
\item An open environment, checkpoints, per-episode logs, an interface-flag
test suite, and a reproduction script that regenerates every table and figure
in this paper from the committed CSVs without a simulator run.
\end{enumerate}

% ------------------------------------------------------------------------
\section{Related work}
\label{sec:related}

\paragraph{Precision limits of learned policies.}
\citet{xu2026curse} introduce the limit precision $c$ and the scaling law
$\log N \propto 1/(P-c)$ for behaviour cloning with diffusion policies; the
action interface is fixed and RL is left to future work.
\citet{gyenes2026fourier} show that MLP-based imitation policies fail to learn
high-frequency components of the state-to-action map and that Fourier feature
encodings \citep{tancik2020fourier} restore precision; we adopt their encoding
as a control arm so that a representation explanation of our floor can be
tested rather than assumed away. High-precision RL results exist. \citet{luo2025hilserl} are the strongest:
their human-in-the-loop system trains directly on a Franka and reaches
near-perfect success on tasks with sub-millimetre clearances, with an average
2$\times$ improvement over imitation baselines, in one to two and a half hours
per task. \citet{brown2026match} reach sub-millimetre-clearance insertion with a
hybrid position/force interface at 15\,Hz. In both cases the interface is a
design contribution rather than a swept variable, and in the first it is
entangled with human intervention and a learned reward classifier. Our question
is complementary and deliberately narrower: holding the human, the reward and
the data collection fixed --- indeed absent --- how much of the terminal
precision does the interface alone account for? A demonstration that a
combination works cannot answer that; only a sweep can.

\paragraph{Action-space design.}
The most closely related work is \citet{aljalbout2024role}, who train over 250
RL agents across 13 action spaces (joint/Cartesian $\times$
position/velocity/torque $\times$ one- and multi-step delta integrators) on
reaching and pushing, at a fixed 60\,Hz, and find that ``an action space which
yields or naturally limits the tracking error transfers better.'' They vary the
\emph{type} of the interface; we hold the type fixed (Cartesian delta position
through differential IK) and vary its \emph{magnitude, filtering and rate}.
\citet{feng2026demystifying} sweep delta/absolute, joint/task, chunk length and
execution horizon for imitation learning and recommend delta actions;
\citet{azimi2026benchmarking} compare pose- and joint-space increments and
velocities for vision-based RL and find joint velocity transfers best. Neither
reports a precision floor or varies step size, smoothing or frequency.
\citet{bronars2026tune} re-tune hyperparameters for each PD-gain regime and find
RL indifferent to gains ``as long as the remaining hyperparameters are tuned to
be compatible''; their frequency ablation (10--100\,Hz) targets jitter in
sim-to-real reaching. Their result is the strongest prior reason to expect
$c(\iface)$ to be flat, and our hyperparameter-optimisation arm
(Sec.~\ref{sec:design}) is designed to test exactly that. Action chunking in RL
\citep{li2025chunking,chen2026acsac} changes the effective control rate and
temporal smoothness of the interface for long-horizon credit assignment; our
frequency and smoothing axes isolate the same two quantities for precision.

\paragraph{Simulator fidelity as a confound.}
\citet{aljalbout2025reality} warn that policies may succeed ``simply by abusing
modeling inaccuracies and simulator-specific corner cases'';
\citet{kaup2024review} find that agents trained in some engines do not transfer
to others and call the same-task-across-engines study missing.
\citet{hu2025static} show a single unmodelled parameter (static friction) can
dominate transfer. Domain randomisation is the standard mitigation
\citep{muratore2022review}, though which parameters matter is usually settled by
prior knowledge \citep{kaidanov2024role}. We do not attempt cross-engine
transfer; we sweep the solver parameters most likely to shape terminal contact
(timestep, constraint \texttt{solref}, impedance ratio, iteration count) and ask
whether $c$ moves. Sec.~\ref{sec:killgate} reports the answer.

\paragraph{Evaluation statistics.}
\citet{jiang2026benchmarking} find that only 19.8\% and 19.7\% of
state-of-the-art claims on LIBERO and SimplerEnv are provably significant,
because aggregate success rates at small $n$ cannot separate methods;
\citet{kressgazit2024empirical} and \citet{snyder2025step} argue for interval
estimates and sequential tests, and \citet{patterson2024empirical} show that
five seeds are almost never enough for strong claims. This is not an abstract
concern for us: an earlier version of this draft reported, from $n{=}200$, that
our two released policies were indistinguishable below a 15\,mm tolerance
($15.0\%$ vs.\ $15.5\%$ at 10\,mm, $p{=}0.89$) and used that as the motivating
observation. At $n{=}2000$ the same comparison gives $17.5\%$ vs.\ $20.8\%$
($z{=}2.65$, $p{=}0.008$): the null was an artefact of the sample size, and the
motivating sentence has been rewritten (Sec.~\ref{sec:baselines}). We therefore
evaluate $n{=}2000$ episodes per cell with Wilson intervals
\citep{wilson1927probable} and three training seeds, and report interquartile
means over seeds \citep{agarwal2021deep}.

\paragraph{Human-in-the-loop RL and its open simulation companion.}
The practical route to \citet{luo2025hilserl}'s result for most groups is
\texttt{gym-hil} \citep{gymhil2025} and the LeRobot port of the method, which
provide MuJoCo Franka tasks with the same capped delta-action interface and the
same human-intervention hooks. We use \texttt{gym-hil} not as a baseline to
beat but as a second measuring instrument: it is the closest public analogue of
our setup written by other people, and repeating one axis of the sweep in it is
the cheapest available guard against a result that is really about the Piper.
Sec.~\ref{sec:twoenvs} reports its released interface constants, which are not
stated in its documentation and which we read out of the code.

\paragraph{What could not be reused.}
We searched for public data that would answer any part of this question without
new rollouts and found none. Manipulation evaluation harnesses discard the
continuous quantity at the moment they threshold it: ManiSkill3's evaluation
returns booleans (the goal distance is computed against a fixed threshold and
dropped), \texttt{robomimic} returns aggregate success rates with no per-episode
dump, and LIBERO exposes only a sparse reward. Demonstration datasets store
object poses and therefore contain terminal errors, but for demonstrations, not
for policy rollouts. \citet{xu2026curse} release nothing;
\citet{feng2026demystifying} promise code on publication and have not yet
released it. Consequently every success-versus-tolerance curve in this paper had
to be generated by re-running policies, and we release ours per-episode so that
the next study does not.

% ------------------------------------------------------------------------
\section{Problem statement}
\label{sec:problem}

\paragraph{Task and interface.}
An episode of the pick-and-place task samples an object pose and a target
position, both drawn from randomised annuli in front of the arm
(Sec.~\ref{sec:setup}). The policy $\pi_\phi$ observes a 51-dimensional state
and outputs $a_t \in [-1,1]^5$: a Cartesian displacement of the tool centre
point, a wrist-yaw increment and a gripper command. The environment maps $a_t$
to a joint target through the interface
\begin{equation}
  \Delta x_t = \dmax\, a_t^{xyz},\qquad
  q^{\mathrm{raw}}_t = \mathrm{IK}(x_{t-1} + \Delta x_t),\qquad
  q^{\mathrm{cmd}}_t = (1-\alpha)\, q^{\mathrm{cmd}}_{t-1} + \alpha\, q^{\mathrm{raw}}_t,\qquad
  \|q^{\mathrm{cmd}}_t - q^{\mathrm{cmd}}_{t-1}\|_\infty \le \qdmax / f,
  \label{eq:interface}
\end{equation}
executed by position actuators at $f$\,Hz. We call $\iface = (\dmax,\alpha,f,\qdmax)$
the action interface. In the released environment
$\iface_0 = (30\,\mathrm{mm},\ 0.45,\ 20\,\mathrm{Hz},\ 1.0\,\mathrm{rad\,s^{-1}})$.
Two properties of Eq.~\ref{eq:interface} are easy to misread and matter for the
design. First, the smoothing acts on \emph{joint targets after IK} --- it is a
first-order low-pass filter on the commanded trajectory, not on the action.
Second, the joint-velocity clamp is part of a safety layer shared with the
hardware interface, and it is not slack: over the $n{=}2000$ baseline
evaluations it is active on 78\% (\(\pi_{v2}\)) and 83\% (\(\pi_{v1}\)) of
control steps. It is therefore a binding constraint at $\iface_0$ and must be
swept jointly with $\dmax$ rather than held fixed (Sec.~\ref{sec:design}).

\paragraph{Precision and the floor.}
The terminal placement error is $e = \|p^{\mathrm{obj}}_{xy} - p^{\mathrm{tgt}}_{xy}\|_2$
once the object has settled after release. For a tolerance $\tau$ the success
rate is $S(\tau) = \Pr[e \le \tau \wedge \text{lifted} \wedge \text{settled}]$;
an episode that never lifted or never settled counts as a failure at every
tolerance, so that a dropped object lying 0.1\,mm from the target is not scored
as a 0.1\,mm success. Because $e$ is logged per episode, a single evaluation
yields $S(\tau)$ for every $\tau$ at once, and
\texttt{tests/test\_tolerance\_semantics.py} asserts that reading the curve post
hoc agrees exactly with thresholding at that $\tau$. We use two estimators of
the floor. The \emph{plateau estimator} $c_{50}$ is the tolerance at which
$S(\tau)$ crosses 50\%; it is model-free, and equals the median terminal error
whenever the valid fraction exceeds one half. The \emph{scaling-law estimator}
$\hat c$ fits Xu et al.'s form to precision-versus-training-steps curves
obtained by evaluating every stored checkpoint of a training run, replacing
demonstrations $N$ by environment steps; it is comparable to their $c$ but
model-dependent. We also report $c_{90}$, the 90th-percentile tolerance, which
is more sensitive to the tail and correspondingly noisier.

\paragraph{Hypotheses.}
\begin{description}[nosep,leftmargin=1.5em]
\item[H1 (interface).] $c$ is a monotone function of the effective per-step
resolution of $\iface$. Under a one-factor-at-a-time sweep, $c_{50}$ moves with
$\dmax$, $\alpha$ and $f$ by more than its 95\% interval.
\item[H0a (reward).] Adding a precision term to the terminal reward moves $c$.
\item[H0b (budget).] Training three times longer at fixed $\iface_0$ moves $c$.
\item[H0c (representation).] Fourier-feature encoding of the observation moves
$c$ at fixed $\iface_0$.
\item[H0d (simulator).] Changing timestep, \texttt{solref}, impedance ratio or
solver iterations within ranges used by published MuJoCo manipulation work
moves $c$ at fixed $\iface_0$.
\item[H0e (release physics).] The floor is dominated by object scatter between
release and settling rather than by the TCP error at the moment of release.
\end{description}
H1 is supported if it holds and H0a--H0e fail. If H0d holds, the paper's claim
collapses to ``the floor is a simulator artefact'' and we say so; this is the
pre-registered kill gate. If H0c holds, the representation explanation of
\citet{gyenes2026fourier} extends to RL and H1 must be restated as a joint
interface--representation account. H0d and H0e are settled in
Sec.~\ref{sec:phase0}; H0a--H0c and H1 require the training campaign.

% ------------------------------------------------------------------------
\section{Experimental setup}
\label{sec:setup}

\paragraph{Environment.}
The arm is an AgileX Piper (6 revolute joints plus a parallel gripper) modelled
in MuJoCo \citep{todorov2012mujoco} from the Menagerie-derived MJCF, with
\texttt{implicitfast} integration, elliptic friction cones, \texttt{impratio}=10
and otherwise default solver settings (timestep 2\,ms, Newton solver, 100
iterations, 25 physics sub-steps per control step at 20\,Hz). Gravity
compensation is enabled on all links. Position actuators use $k_p{=}80,\ k_v{=}5$
on joints 1--3, $k_p{=}40$ on joint 4, $k_p{=}10,\ k_v{=}1.5$ on joints 5--6 and
$k_p{=}200$ on the gripper. The object is a cylinder of radius 24\,mm and height
70\,mm, mass 60\,g, with contact parameters \texttt{solref}=(0.010,\,1.2),
\texttt{solimp}=(0.92,\,0.97,\,0.001) and friction (1.5,\,0.05,\,0.005); finger
pads use \texttt{solref}=(0.012,\,1.4). The target is a site on the table;
success requires the object within $\tau_{xy}$ of the target and 20\,mm in
height, at rest (speed $<0.05$\,m/s), for three consecutive control steps with
the gripper open. The default tolerance is $\tau_{xy}{=}45$\,mm; all curves below
are evaluated post hoc at every $\tau$.

\paragraph{Observation and action.}
The 51-dimensional state contains normalised joint positions and velocities,
gripper state, TCP position and approach vector, object and target positions,
object velocity, relative vectors, grasp flags, the previous action and the
sampled object size. The 5-dimensional action is described in
Eq.~\ref{eq:interface}; differential IK uses damped least squares (damping 0.06,
step 0.6, two iterations) with a position-first null-space projection.

\paragraph{Domain randomisation and noise.}
Per episode: object radial distance 0.30--0.39\,m and azimuth $[-0.73,-0.17]$\,rad,
target on the mirrored annulus; object radius and height scale $[0.85,1.15]$,
mass 30--120\,g, friction 0.8--2.0; table friction 0.6--1.4; actuator $k_p$ and
$k_v$ scales $[0.8,1.25]$ per actuator; joint friction-loss scale $[0.5,1.6]$;
initial joint jitter 0.06\,rad; lighting. Observation noise is additive Gaussian
(joint position 3\,mrad, object position 6\,mm, target 4\,mm, object rotation
0.05\,rad) with one step (50\,ms) of observation latency, 4\,mrad action noise
and 1\% command dropout. A ``hard corner'' fraction oversamples the far-radius
region during fine-tuning. \textbf{Correction:} the environment also draws a
$\pm40\%$ joint-damping scale, and earlier descriptions of this work listed it
among the randomised dynamics. The released MJCF sets no joint damping, so that
scale multiplies zero and is a no-op. The randomiser now reports the drawn
scales together with a \texttt{joint\_damping\_is\_noop} flag, and the flag is
true for the released model; the parameter is excluded from the randomisation
count and from the failure attribution. Adding nominal damping to the MJCF
would change the dynamics of every result in this paper and was therefore not
done after the evaluations began.

\paragraph{Learning algorithm.}
Soft actor-critic \citep{haarnoja2018sac} from Stable-Baselines3
\citep{raffin2021sb3}: two 256-unit hidden layers, learning rate
$3{\times}10^{-4}$ (fine-tuning $10^{-4}$), batch 256, replay buffer
$4{\times}10^5$, $\gamma{=}0.98$, $\tau{=}0.01$, automatic entropy tuning
initialised at 0.1, 12 parallel environments, 12 gradient steps per environment
step, 5\,000 warm-up steps. A reverse curriculum samples episode starts from
later task phases with decreasing probability during training and is disabled at
evaluation. The reward combines dense reaching, lifting and transport shaping
with stage bonuses (touch 3, grasp 25, lift 25, over-target 12, success 200) and
per-step penalties on time, action magnitude, action change, joint speed,
collision and unsafe commands. A terminal precision term
$w_{\mathrm{prec}}(1 - e/\tau_{xy})$ exists and is \emph{disabled}
($w_{\mathrm{prec}}{=}0$) in both released checkpoints. Measured throughput on
the development machine (CPU-only training) is 26--62 environment steps per
second; a 1.2\,M-step run takes 5--6\,h. \todo{Re-measure throughput with SAC
updates on the GPU before the sweep; the sweep budget below assumes $\le 6$\,h
per run.}

\paragraph{The two released policies, and what separates them.}
$\pi_{v1}$ is the \emph{best}-checkpoint selection of a 1.0\,M-step run
(\texttt{runs/sac\_full}); $\pi_{v2}$ is the \emph{final} model of a 0.3\,M-step
fine-tune of it (\texttt{runs/corner\_ft}) with 35\% hard-corner oversampling
and learning rate $10^{-4}$. The reward is identical in both
($w_{\mathrm{prec}}{=}0$); an earlier description of this work stated that they
differed by reward, and that was wrong. Any $\pi_{v1}$-vs-$\pi_{v2}$ difference
therefore confounds three things --- the shifted sampling distribution, the
extra 0.3\,M steps, and the checkpoint-selection rule --- and none of them is
the reward. This is exactly why the sweep in Sec.~\ref{sec:design} retrains
every cell, including the $\iface_0$ cell, under one protocol with final-checkpoint
selection, rather than reusing either released policy as a baseline.

\paragraph{Evaluation protocol.}
Deterministic policy, randomisation and noise on, curriculum off, episode $i$
seeded with $s_0+i$. Every episode writes one CSV row: seed, git hash, success,
grasp, lift, termination reason, terminal error $e$, TCP--target error at the
release step, post-release scatter, episode length, collision steps, safety
clamp count and rate, the resolved interface and solver settings, and the full
sampled randomisation vector. The release step is detected physically, as the
last control step on which both finger pads are in contact with a lifted object,
rather than from the gripper command --- a command-based detector fires on
transient open commands during transport and inflates both terms of the
decomposition. Success curves $S(\tau)$ are reported with 95\% Wilson intervals
\citep{wilson1927probable}; differences between cells are tested with a
two-proportion $z$-test at the tolerance of interest and reported with effect
size in percentage points and a 95\% interval; $c_{50}$ and $c_{90}$ carry a
2000-draw percentile bootstrap over episodes. Each training cell will be run
with three seeds; we report the per-seed curves and the interquartile mean.

\paragraph{Instrumentation as a first-class artefact.}
Two test files guard the sweep. \texttt{tests/test\_interface\_flags.py}
asserts, for every swept quantity, that the flag reaches the simulator: that
$\dmax$ bounds the one-step displacement and is monotone in reachable
displacement, that the control rate changes the number of physics sub-steps
($10/20/50$\,Hz $\to 50/25/10$) \emph{and} the safety clamp $\qdmax/f$, that
$\qdmax$ sets that clamp, that the solver overrides land on
\texttt{model.opt} and on the object and finger-pad \texttt{solref} but not on
the table, and that scaling \texttt{solref} preserves the damping ratio. Without
these, a mis-wired flag would make several cells of Figure~\ref{fig:sweep} the
same experiment run repeatedly, and nothing in the printed output would say so.
One test documents a real defect rather than asserting correctness: at 20\,Hz a
4\,ms timestep is 12.5 sub-steps and silently rounds to 12, a 4\% shortfall in
simulated time per control step. The eval-only gate reports the 4\,ms cell as
measured; the retrained solver extreme uses 2.5\,ms, which divides exactly.

% ------------------------------------------------------------------------
\section{Experimental design}
\label{sec:design}

The design is one-factor-at-a-time around $\iface_0$, with retraining in every
cell. Table~\ref{tab:design} lists the cells. Phases are ordered so that the
experiment most likely to kill the hypothesis runs first, and so that
everything obtainable without training is obtained first.

\begin{table}[t]
\centering\small
\caption{Cells. Every training cell is run with three seeds and evaluated with
$n{=}2000$ episodes. ``Eval-only'' cells re-evaluate an existing checkpoint
under a changed simulator or tolerance and cost no training. Run-time estimates
assume 6\,h per Piper training run; the \texttt{gym-hil} cells are roughly an
hour each. Phase~0 is complete; Phases~1--2 are pending.}
\label{tab:design}
\resizebox{\linewidth}{!}{\begin{tabular}{@{}llllrr@{}}
\toprule
Phase & Arm & Cells & Purpose & Runs & Status \\
\midrule
0 & Solver (kill gate), eval-only & timestep $\{1,2,4\}$\,ms; \texttt{solref} scale $\{0.5,1,2\}$; iterations $\{20,100\}$; \texttt{impratio} $\{1,10\}$; two joint extremes & does the existing floor move? & 0 & done \\
0 & Checkpoint curve, eval-only & all 38 stored checkpoints of both runs & $\hat c$ vs.\ training steps; H0b & 0 & done \\
0 & Release decomposition, eval-only & $\pi_{v1},\pi_{v2}$ & H0e & 0 & done \\
0 & Nested noise replay, eval-only & 200 scenes $\times$ 10 noise streams, $\pi_{v2}$ & ICC; failure attribution & 0 & done \\
0 & Solver (kill gate), retrain & two extreme settings & H0d & 6 & pending \\
\midrule
1 & $\dmax$ & $\{5, 10, 20, 30\}$\,mm, with $\qdmax$ scaled so the clamp binds equally & H1 & 12 & pending \\
1 & $\alpha$ & $\{0.2, 0.45, 1.0\}$ (1.0 = no smoothing) & H1 & 6 & pending \\
1 & $f$ & $\{10, 50\}$\,Hz & H1 & 6 & pending \\
1 & $\qdmax$ alone & $\{0.5, 2.0\}$\,rad/s at $\dmax{=}30$\,mm & disentangle clamp from step & 6 & pending \\
1 & \textbf{2nd env.} \texttt{gym-hil} & $\dmax \in \{5,10,25,50\}$\,mm, $\alpha \in \{0.2,0.45\}$, $f \in \{5,25\}$\,Hz & is H1 about interfaces or about the Piper? & 24 & pending \\
\midrule
2 & Reward & $w_{\mathrm{prec}}{=}60$ from scratch & H0a & 3 & pending \\
2 & Budget & $3{\times}$ steps at $\iface_0$ & H0b & 1 & pending \\
2 & Fourier features & $\iface_0$ with Fourier-encoded state & H0c & 3 & pending \\
2 & HPO & 10 Optuna trials at $\dmax{=}5$\,mm & \citet{bronars2026tune} objection & 10 & pending \\
\midrule
 & & & Total training runs & 77 & \\
\bottomrule
\end{tabular}}
\end{table}

\paragraph{Phase 0: kill gate and free measurements.}
Before any training, the released policy $\pi_{v2}$ is re-evaluated with each
solver setting changed in the loaded model (timestep and iterations through
\texttt{model.opt}, contact parameters by scaling \texttt{geom\_solref} on the
object and the four finger-pad geoms, preserving the damping ratio). If the
released policy's $c_{50}$ moves by more than its interval under a solver
change, we retrain at the two extreme settings; if the retrained floor moves as
far as it does under a $30\to10$\,mm step change, the study stops here and the
write-up becomes a negative result about simulator dependence. The same phase
evaluates every stored checkpoint of both runs to obtain precision as a function
of training steps, and decomposes each terminal error into TCP-at-release error
and post-release scatter. Sec.~\ref{sec:phase0} reports all three.

\paragraph{Phase 1: interface sweep.}
Cells in $\dmax$ are run with $\qdmax$ scaled proportionally
($30\,\mathrm{mm}\to1.0$, $20\to0.67$, $10\to0.33$, $5\to0.17$\,rad/s), because
at $\iface_0$ the joint-velocity clamp binds on about four steps in five
(Sec.~\ref{sec:baselines}) and a step-size sweep with a fixed clamp would
measure the clamp. A separate two-cell sweep of $\qdmax$ at fixed $\dmax$
separates the two, and the clamp rate is reported per cell. When $f$ changes,
the number of physics sub-steps, the observation latency in milliseconds and the
episode budget in steps all change; we hold latency and budget fixed in seconds,
so only the sub-step count and the clamp --- both parts of the interface --- move
with $f$. At 50\,Hz a 30\,mm step is $1.5\,\mathrm{m\,s^{-1}}$, above the
safety layer's Cartesian rate limit of $0.6\,\mathrm{m\,s^{-1}}$; that cell
raises the limit to $1.6\,\mathrm{m\,s^{-1}}$ so that a second, unswept limiter
does not silently bind, and the change is reported. Domain randomisation ranges
are held fixed at their released values in all cells, since
\citet{xu2026curse} show $c$ moves with randomisation range.

\paragraph{Phase 2: controls.}
The reward arm trains from scratch with the precision term enabled, rather than
fine-tuning, so that budget and reward are not confounded. The budget arm
continues the baseline to $3.6$\,M steps. The Fourier arm replaces the state
input by $[s, \sin(2\pi B s), \cos(2\pi B s)]$ with a fixed Gaussian $B$
following \citet{gyenes2026fourier}, using the same $B$ in every parallel
environment and at evaluation. The HPO arm re-tunes learning rate, entropy
coefficient and batch size at the smallest step size using Optuna
\citep{akiba2019optuna}; if tuned hyperparameters recover the $\iface_0$ floor at
$\dmax{=}5$\,mm, the Bronars-style objection stands and we report it, and the
claim weakens to ``the interface sets the floor for a fixed optimiser
configuration.''

\paragraph{Failure attribution.}
\label{sec:attrib}
For every cell we fit a cross-validated logistic model and a gradient-boosted
classifier of failure at $\tau{=}c_{50}$ on the sampled randomisation vector and
initial state, and report AUC and permutation importances. A nested design (1000
fixed randomisation draws $\times$ 10 noise re-seeds) estimates the intra-class
correlation of outcomes within a draw, separating environment-determined from
irreducible failures. \planned{If AUC $<0.6$ and ICC $<0.1$ in every cell, the
failures are noise and this section reduces to a paragraph.}

% ------------------------------------------------------------------------
\section{Results}
\label{sec:results}

\subsection{Phase 0, measured: what the released checkpoints already say}
\label{sec:phase0}

Everything in this subsection was obtained without training a single policy. It
uses only the checkpoints committed in the repository, and it is reproduced end
to end by \texttt{python reproduce.py} from the archived per-episode CSVs.

\subsubsection{The floor of the two released policies}
\label{sec:baselines}

\input{generated/tables/baselines}

Table~\ref{tab:baselines} gives the full success-versus-tolerance curve for both
released policies at $n{=}2000$ episodes on a common seed block, with 95\%
Wilson intervals, and $c_{50}$ with a 2000-draw bootstrap. Four things follow.

\textbf{The cliff is steep and it is not near zero.} $\pi_{v2}$ succeeds on
94.9\% of episodes at $\tau{=}45$\,mm and on 5.3\% at $\tau{=}5$\,mm; its
$c_{50}$ is $16.1$ [15.7, 16.6]\,mm and its $c_{90}$ is $28.9$\,mm. The floor
sits at roughly half the 30\,mm maximum step and at 1.4 times the effective
per-step reach of the joint-velocity clamp, both of which are consistent with
H1 and neither of which is evidence for it: they are the quantities the sweep
must separate.

\textbf{The joint-velocity clamp is the binding interface term at $\iface_0$.}
It fires on 78\% of the control steps of an average $\pi_{v2}$ episode and 83\%
of a $\pi_{v1}$ episode. A step-size sweep that held $\qdmax$ fixed would
therefore be measuring the clamp for most of its range, which is why the
$\dmax$ cells scale $\qdmax$ with them and a separate axis varies $\qdmax$
alone.

\textbf{The motivating null of the earlier draft does not survive $n{=}2000$.}
An earlier version of this work reported that the two policies were
indistinguishable below 15\,mm ($15.0\%$ vs.\ $15.5\%$ at $\tau{=}10$\,mm,
$n{=}200$, $p{=}0.89$) and used that as the observation motivating the whole
study. At $n{=}2000$ the difference at 10\,mm is $+3.3$ percentage points
[$+0.9$, $+5.7$], $z{=}2.65$, $p{=}0.008$, and it is still significant at
7.5\,mm ($+2.6$\,pp, $p{=}0.007$); it disappears only at 5\,mm ($+0.9$\,pp,
$p{=}0.21$). The two $c_{50}$ intervals, $[18.1, 19.4]$ and $[15.7, 16.6]$\,mm,
do not overlap. Training history \emph{does} move this floor. What survives is
the weaker and more interesting statement that the \emph{effect} of training
history decays monotonically as the tolerance tightens --- $+13.5$\,pp at
20\,mm, $+8.5$ at 15, $+3.3$ at 10, $+2.6$ at 7.5, $+0.9$ at 5 --- so the two
histories converge as the ruler approaches the floor without ever having been
equal. We report this because it is the clearest illustration in the paper of
the point \citet{jiang2026benchmarking} and \citet{snyder2025step} make: at
$n{=}200$ we would have published a null that is not there.

\textbf{The comparison is confounded anyway.} $\pi_{v1}$ and $\pi_{v2}$ differ
by a shifted sampling distribution, 0.3\,M extra steps \emph{and} the
checkpoint-selection rule (best vs.\ final). No causal reading of the difference
is available, which is precisely why every cell of Phase~1 is retrained under a
single protocol.

\subsubsection{H0e: the floor is set at release, not by post-release settling}

\input{generated/tables/release}

If the terminal error were dominated by how far the object rolls after the
fingers open, no change to the action interface could move it and H1 would be
irrelevant. Table~\ref{tab:release} decomposes the error. For $\pi_{v2}$ the
median tool-centre-point error at the instant the object leaves the pads is
17.7\,mm (p90 27.9), the median distance the object subsequently travels before
settling is 7.6\,mm (p90 17.0), and the median terminal error is 15.7\,mm. The
terminal error correlates with the TCP error at release ($r{=}0.35$ for
$\pi_{v2}$, $0.53$ for $\pi_{v1}$) and is uncorrelated with the post-release
scatter ($r{=}0.07$ and $-0.06$; $r^2 \approx 0.004$). \textbf{H0e is rejected:}
the object does move after release, but where it ends up is set by where the
gripper let go, not by the rolling. Terminal precision in this task is a
placement problem at the last held step, which is the step the action interface
governs.

\subsubsection{H0d: the contact-solver kill gate}
\label{sec:killgate}

\input{generated/tables/killgate}

Table~\ref{tab:killgate} re-evaluates $\pi_{v2}$ under nine contact-solver
settings spanning the range used by published MuJoCo manipulation work:
integration timestep 1, 2 and 4\,ms; contact time constant of the object and
finger pads scaled by $0.5$, $1$ and $2$ with the damping ratio held; Newton
iterations 20 and 100; impedance ratio 1 and 10; and the two joint extremes
(1\,ms with \texttt{solref}$\times0.5$, 4\,ms with \texttt{solref}$\times2$).
Each cell is $n{=}1000$ on a common seed block.

$c_{50}$ moves by at most 1.9\,mm, a total span of 2.7\,mm across the nine
cells against a nominal 16.4 [15.7, 17.0]\,mm --- about 17\% of the floor. The
movement is not uniformly negligible: two cells, \texttt{solref}$\times0.5$
($18.3$ [17.5, 19.1]\,mm) and the stiff/fine extreme ($17.1$ [16.5, 18.0]\,mm),
have intervals that exclude the nominal point estimate. Softening the contact
time constant makes the floor worse and, more sharply, degrades the tail:
$S(45)$ falls from 95.8\% to 88.6\% and $c_{90}$ rises correspondingly, while
$c_{50}$ moves by under 2\,mm. In other words the solver mostly changes how
often the policy fails outright, not how precisely it places when it succeeds.

Two negative details are worth stating because they would otherwise look like
bugs. The 20-iteration cell is \emph{bit-identical} to the 100-iteration cell in
all 1000 episodes; a direct probe (\texttt{results/precision/solver\_iterations\_probe.md})
shows the knob is live --- 5 iterations changes the trajectory, 2 or 1 breaks
the simulation entirely, with median errors above 450\,mm --- and that the
Newton solver simply converges in fewer than 20 iterations for this contact
configuration. Iteration count in $[20, 100]$ is an identity here, not a
measured null. Separately, at 20\,Hz a 4\,ms timestep is 12.5 sub-steps and
rounds to 12; that cell is reported as measured and the retrained extreme uses
2.5\,ms instead.

\paragraph{Gate decision.} The pre-registered rule was: if $c_{50}$ moves by
more than its interval under any eval-only solver cell, retrain at the two
extreme settings before proceeding, and abandon the interface claim if the
retrained floor moves as far as a $30\to10$\,mm step change does. Two cells
moved by more than their interval, so \textbf{the gate escalates to the
retraining stage}: six runs (two extremes $\times$ three seeds) that have not
been performed. The eval-only bound is nonetheless informative --- a 1.9\,mm
solver effect is the number the interface effect must beat, and if the $\dmax$
axis turns out to move $c_{50}$ by, say, 2\,mm, the paper has no result
regardless of what the retrained extremes do. \todo{Run the six retraining runs
and fill this decision. Do not run Phase~1 before it.}

\subsubsection{H0b, without training anything: precision versus training budget}
\label{sec:ckpt}

\begin{figure}[t]
\centering
\includegraphics[width=\linewidth]{generated/figures/checkpoints}
\caption{Every stored checkpoint of both runs, $n{=}1000$ episodes each (38
checkpoints, 38\,000 episodes, no training). Left: success at the released
45\,mm tolerance. Right: $c_{50}$, which is undefined ($\times$) until the
policy succeeds on more than half of episodes; shaded bands are bootstrap 95\%
intervals. The floor is a property of the converged policy: it is undefined for
the first 0.82\,M steps, drops sharply once the task is solved, and then does
not improve.}
\label{fig:ckpt}
\end{figure}

Figure~\ref{fig:ckpt} evaluates all 38 stored checkpoints --- 32 of the baseline
run spanning 50\,k to 1.62\,M environment steps and 6 of the fine-tune --- at
$n{=}1000$ each. It is the RL analogue of Xu et al.'s data-scaling curve with
environment steps in place of demonstrations, and it costs no training at all.

\textbf{$c_{50}$ does not improve with budget once the task is solved.} For the
baseline run it is undefined below 0.82\,M steps (success never reaches 50\%),
first measurable at 31.4\,mm at 0.87\,M, falls to $16.8$\,mm at 1.22\,M, and
then \emph{oscillates between 16.8 and 24.3\,mm} for the remaining 0.4\,M steps,
finishing at $24.3$ [23.7, 25.2]\,mm at 1.62\,M with success at 45\,mm falling
from 96.5\% to 74.3\%. Over the last third of the run the floor moves by more
than its own bootstrap interval in both directions and ends \emph{worse} than
its best. This is a within-run, single-seed observation and it does not replace
the pre-registered $3\times$-budget arm, which is trained from scratch --- but
it is the first evidence on H0b and it points the same way: past the point where
the task is solved, more environment steps buy noise, not precision.

\textbf{Two consequences for how the released policies should be read.} First,
$\pi_{v2}$ ($c_{50}{=}16.1$\,mm) is not the best checkpoint of its own
fine-tune: the checkpoint at 1.06\,M steps reaches $14.6$ [14.0, 15.3]\,mm.
Second, the first fine-tune checkpoint, immediately after the baseline weights
are loaded and the replay buffer is empty, sits at $35.9$\,mm --- a large
transient dip that recovers within 50\,k steps. A campaign that selected
checkpoints by a 15--50-episode evaluation callback, as the released runs did,
is selecting from this much noise, which is why the sweep in
Sec.~\ref{sec:design} fixes the final checkpoint as the reported one.

\textbf{Xu et al.'s functional form does not fit.} Fitting
$\log(\text{steps}) = a/(P - c) + b$ with $P = 1/c_{50}$ over the checkpoints
that solve the task gives $\hat c = 18.4$\,mm with $r^2 = 0.17$ for the baseline
run and an uninterpretable fit ($r^2 \approx 0$, six points) for the fine-tune.
We report the failure rather than the number. A behaviour-cloning data-scaling
curve is monotone by construction --- more demonstrations never hurt --- whereas
an RL checkpoint curve is a stochastic process that can and does go backwards,
so the model-dependent estimator $\hat c$ is not usable in this setting and the
model-free $c_{50}$ is the estimator we report throughout. That is itself a
finding about transporting their result to RL, and it is one we would have
missed by reporting a fitted constant without its $r^2$.

\subsection{The second environment, and what its released interface is}
\label{sec:twoenvs}

The step-size axis is repeated in \texttt{gym-hil} \citep{gymhil2025}, the
MuJoCo suite that accompanies the human-in-the-loop line of work. Its Franka
Panda is driven by operational-space control towards a mocap target, not by
differential IK with a joint-velocity safety clamp; it was written by different
people for a different robot; and it is Apache-2.0, so the comparison is
reproducible by anyone.

\begin{table}[t]
\centering\small
\caption{The two environments' released action interfaces, read out of the code
rather than from documentation. Neither project states these numbers in its
README, and in \texttt{gym-hil} the base class's signature (50\,Hz) disagrees
with the task environment's default (10\,Hz). The last row is the observation
that motivates H1: three tasks written by two groups for two arms all set their
success tolerance between one and two commanded step sizes.}
\label{tab:twoenvs}
\begin{tabular}{@{}lcc@{}}
\toprule
 & this paper (Piper) & \texttt{gym-hil} (Franka) \\
\midrule
controller & differential IK + safety layer & operational-space, mocap target \\
$\dmax$ (commanded) & 30\,mm & 25\,mm \\
$\dmax$ (realised in one step) & --- & $\approx$10\,mm ($0.40\times$) \\
$f$ & 20\,Hz & 10\,Hz \\
physics sub-steps & 25 $\times$ 2\,ms & 50 $\times$ 2\,ms \\
$\alpha$ & 0.45 & 1.0 (none) \\
$\qdmax$ & 1.0\,rad/s, binds on 78--83\% of steps & none \\
episode budget & 200 steps = 10\,s & 100 steps = 10\,s \\
success tolerance & 45\,mm (placement) & 50\,mm (grasp), 30\,mm (arrangement) \\
\midrule
tolerance $/\ \dmax$ & 1.5 & 2.0, 1.2 \\
\bottomrule
\end{tabular}
\end{table}

Table~\ref{tab:twoenvs} reports the constants. Two of them are worth
calling out because they are not documented anywhere and a user of the package
would not know them. First, the 25\,mm cap lives in \texttt{EEActionWrapper}
and applies only to the wrapped environment ids; the \texttt{Base} ids apply no
scaling at all, so an action of $1.0$ commands a one-metre mocap displacement,
clipped to the workspace box. Second, the \emph{realised} displacement of a
saturated step is about 40\% of the commanded cap --- roughly 10\,mm for the
released 25\,mm setting --- because operational-space control does not reach the
mocap target within one 100\,ms control step. Any comparison of the two
environments' step sizes has to use the realised number; the ratio is linear in
$\dmax$ and is pinned by a test
(\texttt{tests/test\_gymhil\_interface.py}) so that a future release of the
package cannot change it silently.

\paragraph{Which task, and an honest reason.}
The natural analogue of our placement task is \texttt{PandaArrangeBoxes},
which has a real 30\,mm placement tolerance. As released it is not learnable
from state: success is a conjunction over five blocks, while the state
observation is three numbers --- the position of block~1. The other four blocks
and all five targets are unobservable. We therefore sweep
\texttt{PandaPickCube} and take the terminal gripper-to-block distance as the
precision measure, which is exactly the quantity \texttt{gym-hil}'s own success
condition thresholds at 50\,mm. That is a reach-and-grasp precision rather than
a placement precision, and the difference is a limitation of the comparison
rather than a feature of it. A test records the defect so that the task choice
is revisited if a future release fixes it.

\planned{Results: $S(\tau)$ and $c_{50}$ for $\dmax \in \{5, 10, 25, 50\}$\,mm
on \texttt{PandaPickCube}, three seeds, $n{=}2000$ each, with the $\alpha$ and
$f$ axes as secondary cells. If $c_{50}$ tracks $\dmax$ in both environments,
H1 is a statement about delta-action interfaces; if it tracks it in one, H1 is
a statement about the Piper and the paper says so.} These cells are cheap ---
throughput is 63 environment steps per second on two CPU cores against 26--62
for the Piper with 12 gradient steps, and the budget is 300\,k steps rather
than 1.2\,M, so a three-seed step axis is under a day of wall clock.
\todo{Run and fill; 24 runs via \texttt{run\_sweep.py --env gymhil}.}

\subsection{Phase 1: interface sweep}
\planned{Figure 1 (the paper's main figure): $S(\tau)$ curves per cell with
intervals, one panel per axis. Table: $c_{50}$ and $c_{90}$ per cell, IQM over
seeds. Regression of $c_{50}$ on the effective per-step resolution
$\min(\dmax, \qdmax\,J/f)\cdot g(\alpha)$.}
\todo{Run and fill; 30 training runs. Blocked on the gate decision above.}
\label{fig:sweep}

\subsection{Phase 2: controls}
\planned{One table: $c_{50}$ at $\iface_0$ for baseline, reward arm, budget arm,
Fourier arm; and $c_{50}$ at $\dmax{=}5$\,mm with default vs.\ tuned
hyperparameters.}
\todo{Run and fill; 17 training runs.}

\subsection{Failure attribution: how much of the residual failure is the scene?}
\label{sec:attrib-results}

At $\tau = c_{50}$, by construction, half the episodes fail. Two questions
follow: can the failures be predicted from the sampled randomisation vector, and
would the same physical scene fail again?

\textbf{Prediction.} A five-fold cross-validated classifier of failure on the
thirteen sampled randomisation and scene variables reaches AUC $0.56$--$0.65$
across the Phase~0 cells at $\tau = c_{50}$ and $0.53$--$0.65$ at 45\,mm ---
better than chance, but only just. Permutation importance puts the target's
azimuth and radius first in almost every cell, with the object's sampled width
next; the object's mass, its friction, the table friction and the actuator gain
scales sit at or below zero. Where the object has to be \emph{put} in the
workspace carries the signal; what it is made of does not.

\textbf{Replay.} A nested evaluation of $\pi_{v2}$ --- 200 randomisation draws,
each replayed under 10 independent noise streams, 2000 episodes --- separates
the two sources of variance directly. The environment supports the split
(\texttt{options=\{"dr\_seed": k, "noise\_seed": m\}}), and the classifier for
this cell is cross-validated with folds grouped by scene, without which replays
of one scene appear on both sides of the split and the AUC is inflated by about
0.08. The one-way random-effects intra-class correlation of the binary outcome
is $\mathrm{ICC} = 0.175$ at $\tau = c_{50}$ and $0.113$ at 45\,mm.

\textbf{Verdict against the pre-specified rule.} The rule was that if
AUC $<0.6$ and ICC $<0.1$ in every cell, the residual failures are noise and
this section collapses to a paragraph. Neither threshold is quite met: about a
sixth of the outcome variance at the floor is attributable to the scene, and a
classifier can see a little of it. The honest statement is therefore the
intermediate one --- \emph{at the floor, roughly 82\% of whether a given episode
succeeds is decided by observation and actuation noise rather than by the scene
it was given}, and the minority that the scene decides is geometric rather than
dynamical. For H1 this is supporting context rather than evidence: a floor that
is five parts noise to one part scene is not a floor that better coverage of the
randomisation space would remove.

% ------------------------------------------------------------------------
\section{Discussion}
\label{sec:discussion}

\todo{Write in full after Phases 1--2. The following is what the completed
Phase 0 already licenses.}

\paragraph{Where the floor is located in the episode.} The decomposition in
Table~\ref{tab:release} narrows the question considerably. The terminal error is
essentially the tool-centre-point error at the last held step; post-release
physics contributes real displacement (a 7.6\,mm median) but almost no
explanatory variance ($r^2 \approx 0.004$). Any account of the floor that
appeals to settling, rolling or contact release is therefore the wrong account
for this task. That is a necessary condition for H1, not a sufficient one: it
localises the floor to the terminal approach, which is where $\dmax$, $\alpha$,
$f$ and $\qdmax$ act --- but also where the observation noise, the IK residual
and the policy's own value estimate act.

\paragraph{How much room the interface hypothesis has.} The kill gate bounds the
simulator's contribution at 1.9\,mm of $c_{50}$ across settings that a
practitioner would consider reasonable. For H1 to be interesting, the interface
sweep must move $c_{50}$ by substantially more than that. It is worth saying now
that the design has no guaranteed win: $\dmax$ could move the floor by 2\,mm and
the paper would then be a null result about interfaces with an awkward
simulator-sensitivity caveat attached. The honest framing in that case is a
short negative note, and it is the framing this draft is prepared to adopt.

\paragraph{Relation to \citet{xu2026curse}.} Their $c$ for peg insertion and
cuboid stacking is 2--3\,mm; ours is 16\,mm on a task with a post-hoc tolerance,
a 30\,mm action step and a low-cost 6-DoF arm. The numbers are not comparable
and we do not compare them. What is comparable is the shape of the claim ---
that $c$ is a property of the agent system rather than the task --- and the
mechanism we propose is a concrete instance of it that they could not observe,
because they held the action interface fixed throughout.

\paragraph{Relation to \citet{bronars2026tune} and \citet{gyenes2026fourier}.}
Both remain open until Phase 2. If hyperparameter re-tuning recovers precision
at small $\dmax$, the interface sets the floor only for a fixed optimiser
configuration, which is weaker but still the practically relevant statement for
anyone who is not going to re-tune. If Fourier features move the floor at fixed
$\iface$, the representation explanation extends to RL with state observations
and H1 must be restated jointly.

\paragraph{Practical reading.} \planned{What a practitioner should change first
when a policy plateaus above the required tolerance --- the interface, the
reward, or the budget --- is the paper's intended takeaway and cannot be written
before Phase 1.} What Phase 0 alone supports is narrower and still useful:
before blaming the simulator, check where in the episode the error is created,
and do not diagnose a floor from a 200-episode evaluation.

% ------------------------------------------------------------------------
\section{Limitations}
\label{sec:limitations}

All experiments are in simulation. The step-size axis is measured on two arms,
two controllers and two tasks --- an AgileX Piper under differential IK with a
joint-velocity safety layer, and a Franka Panda under operational-space control
with none --- which bounds but does not remove the objection that the result is
a property of one robot; the remaining axes and every control arm are measured
on one arm, one task and one algorithm. The second environment's task is a
reach-and-grasp rather than a placement, for the reason given in
Sec.~\ref{sec:twoenvs}, so the two floors are not the same physical quantity
and are compared only in how they move with $\dmax$. The solver arm bounds the
dependence of $c$ on the contact parameters we sweep, not on the contact model
as such; it also found a non-zero dependence, which we report rather than round
to nothing. A real-arm measurement remains the only complete answer, and we do
not have one: \citet{luo2025hilserl} obtain their precision on hardware, and
whether the interface accounts for it there is a question only a lab with that
hardware can settle. What this paper offers is the controlled version --- two
simulators, with the human, the reward classifier and the demonstration data
removed, so the interface can be varied on its own. \todo{If the two-servo encoder rig is
built, report it here as a partial hardware check on the step-size axis only,
with the fixture-compliance calibration; otherwise state plainly that no
hardware validation was performed.} The observation is a ground-truth state;
policies with visual observations may have an additional,
representation-dominated floor. The interface family is Cartesian delta position
through differential IK; joint-space and velocity interfaces
\citep{aljalbout2024role,azimi2026benchmarking} may behave differently. The
task's tolerance is a post-hoc ruler rather than a physical constraint;
\planned{a stacking or insertion variant with a physical tolerance --- a
pedestal of radius $r_{\mathrm{ped}}$ slightly larger than the object, so that
$\tau = r_{\mathrm{ped}} - r_{\mathrm{obj}}$ is a property of the scene ---
would remove this objection and is the natural second task.} Finally, the
released MJCF has no joint damping, so the damping randomisation documented in
earlier versions of this work was inert; the results here are for a model with
zero joint damping, which is itself a modelling choice we did not make
deliberately.

% ------------------------------------------------------------------------
\section*{Reproducibility statement}
The environment, MJCF, training and evaluation scripts, all checkpoints and
every per-episode evaluation log are released at
\url{https://github.com/abdulahad2008/piper_simulation}. \todo{Add a persistent
DOI (Zenodo) and an anonymised mirror before submission.} Every number in
Sec.~\ref{sec:phase0} is regenerated from the committed CSVs by
\begin{quote}\small\texttt{python reproduce.py --results results/precision --out paper/generated}\end{quote}
which reads no simulator and needs no GPU. Package versions are pinned in
\texttt{requirements.txt} and additionally recorded, together with the git hash,
in the \texttt{*.meta.json} beside every results CSV and in every run's
\texttt{config.json}. \texttt{tests/test\_interface\_flags.py},
\texttt{tests/test\_tolerance\_semantics.py} and
\texttt{tests/test\_gymhil\_interface.py} assert that every swept flag reaches
both simulators, that the tolerance means what this paper says it means, and
that \texttt{gym-hil}'s released interface constants are the ones
Table~\ref{tab:twoenvs} reports (75 tests in total). \texttt{CHANGELOG.md} documents every claim that changed between the
August 2026 status documents and the archived numbers, including two that were
simply wrong.

\bibliography{refs}
\bibliographystyle{tmlr}

\appendix
\section{Reward specification}
\todo{Table of every reward term, weight and trigger condition, generated from
\texttt{piper\_rl/config.py} at the commit used.}

\section{Domain randomisation specification}
\todo{Table of every randomised parameter and range, generated from
\texttt{config.py}, including the joint-damping no-op correction and the
\texttt{joint\_damping\_is\_noop} flag now emitted by the randomiser.}

\section{Per-cell training curves}
\todo{Learning curves for every cell and seed.}

\end{document}
